%!TEX TS-program = pdflatex
%!TEX encoding = UTF-8 Unicode

\PassOptionsToPackage{italian,english}{babel}
\documentclass{unipd-thesis-modern}

\usepackage{array}

\setlength{\cftbeforechapskip}{7pt}
\setcounter{tocdepth}{2}

\newcommand{\placeholder}[1]{\textit{[#1]}}
\newcommand{\builder}{buildER}

\title{Relazione finale di tirocinio:\\ sviluppo e consolidamento di \builder}
\thesisHeader{Universita di Bologna\\ \placeholder{Corso di laurea da completare}}
\academicYear{\placeholder{Anno accademico da completare}}

\author{\placeholder{Nome e cognome studente}}
\studentId{\placeholder{Matricola da completare}}

\advisor{Prof.ssa Alessandra Lumini}{Universita di Bologna}

\secondpage{
\mbox{}
\vfill
{\small\onehalfspacing
\begin{center}
\noindent\rule{\textwidth}{1pt}

Relazione finale di tirocinio sul progetto \builder.\\
Tutor del tirocinio: \textbf{Prof.ssa Alessandra Lumini}.\\
Periodo del tirocinio: \placeholder{Periodo del tirocinio da completare}.\\
Ore/CFU: \placeholder{Ore/CFU da completare}.\\
Struttura ospitante: \placeholder{Struttura ospitante da completare}.

\noindent\rule{\textwidth}{.2pt}
\end{center}
}
}

\begin{document}
\selectlanguage{italian}

\makecover
\cleardoublepage

{\hypersetup{linkcolor=black}
\tableofcontents}
\cleardoublepage

\mainmatter

\chapter{Introduzione}

Il presente elaborato descrive il lavoro svolto durante il tirocinio universitario dedicato al progetto \builder, un'applicazione web per la progettazione, visualizzazione, traduzione ed esportazione di schemi Entity-Relationship e schemi logici. Il tirocinio è stato seguito dalla Prof.ssa Alessandra Lumini, tutor del percorso, che ha supervisionato il lavoro e fornito indicazioni utili per orientare le scelte progettuali.

Il progetto si colloca nell'ambito della progettazione di basi di dati, un tema centrale nei percorsi universitari di informatica. La costruzione di uno schema concettuale richiede infatti la capacità di individuare entità, relazioni, attributi, cardinalità e vincoli; la successiva trasformazione verso uno schema logico richiede invece regole coerenti per rappresentare tabelle, colonne, chiavi primarie, chiavi esterne e vincoli relazionali. \builder{} nasce per supportare questo processo in modo interattivo, offrendo un ambiente grafico accessibile da browser.

La relazione non ha l'obiettivo di essere un manuale utente dell'applicazione, ma una descrizione ragionata del percorso svolto durante il tirocinio. Per questo motivo vengono presentati il contesto del progetto, gli obiettivi, le tecnologie utilizzate, l'architettura del sistema, le funzionalità sviluppate o consolidate, le attività di test e i miglioramenti introdotti. Una parte specifica è dedicata anche al ruolo della tutor e alle competenze tecniche e metodologiche acquisite.

L'esperienza ha richiesto di collegare aspetti tecnici e formativi. Da un lato, è stato necessario analizzare un'applicazione web articolata, basata su React, TypeScript e canvas SVG; dall'altro, il lavoro ha richiesto attenzione alla chiarezza degli schemi, alla correttezza concettuale, alla leggibilità dell'interfaccia e alla manutenibilità del codice. Questi aspetti sono stati centrali nel tirocinio e rappresentano il filo conduttore della relazione.

\chapter{Obiettivi del tirocinio}

L'obiettivo generale del tirocinio è stato contribuire allo sviluppo e al consolidamento di \builder{} come applicazione web per la modellazione di schemi ER e per il supporto alla loro trasformazione in schemi logici. Il progetto richiedeva di mantenere un equilibrio tra correttezza concettuale, qualità tecnica e usabilità, poiché uno strumento didattico deve essere comprensibile, affidabile e coerente con i concetti che intende rappresentare.

Gli obiettivi specifici hanno riguardato il miglioramento della modellazione ER, il supporto alla traduzione ER-logico, la leggibilità grafica degli schemi, l'esportazione dei risultati, il salvataggio dei progetti, la localizzazione dell'interfaccia e la verifica tramite test. Tali obiettivi sono stati affrontati considerando sia l'esperienza dell'utente finale sia la struttura interna del codice.

\noindent\begin{tabularx}{\textwidth}{>{\bfseries}p{4cm}X}
Obiettivo & Descrizione \\
\hline
Correttezza concettuale & Rappresentare in modo coerente entità, relazioni, attributi, cardinalità, identificatori e generalizzazioni. \\
Traduzione ER-logico & Supportare il passaggio da schema concettuale a schema logico con tabelle, colonne, chiavi e vincoli. \\
Leggibilità grafica & Migliorare disposizione degli elementi, label, collegamenti e resa degli schemi esportati. \\
Usabilità & Rendere l'interfaccia chiara tramite toolbar, pannelli, modali, scorciatoie e flussi guidati. \\
Esportazione & Produrre output PNG, SVG e JPEG utilizzabili in relazioni, presentazioni e materiali didattici. \\
Persistenza & Salvare e caricare progetti tramite formato \texttt{.ersp}, mantenendo compatibilità con versioni precedenti. \\
Manutenibilità & Organizzare il codice in moduli chiari, separando interfaccia, tipi, logica di dominio e test. \\
Stabilità & Verificare le parti critiche tramite test unitari, test di integrazione e test end-to-end. \\
Valore didattico & Rendere lo strumento utile per comprendere il rapporto tra modello concettuale e modello logico. \\
\end{tabularx}

Dal punto di vista formativo, il tirocinio ha permesso di lavorare su un progetto reale, nel quale ogni scelta tecnica aveva conseguenze sulla comprensibilità dello strumento. La necessità di considerare insieme funzionalità, interfaccia, test e documentazione ha reso il lavoro particolarmente utile per consolidare un metodo di sviluppo più consapevole.

\chapter{Il progetto \builder}

\builder{} è un'applicazione web single-page per la modellazione di schemi Entity-Relationship in stile Chen. L'utente può costruire diagrammi ER su un canvas SVG, aggiungendo entità, relazioni, attributi, cardinalità, identificatori e generalizzazioni ISA. L'applicazione consente inoltre di passare alla vista logica, nella quale lo schema viene rappresentato attraverso tabelle, colonne, chiavi primarie, chiavi esterne e vincoli.

Il progetto è destinato principalmente a un contesto didattico e universitario. Può essere utile agli studenti che devono esercitarsi nella progettazione di basi di dati e ai docenti che vogliono mostrare esempi, discutere schemi o evidenziare il passaggio dal modello concettuale al modello logico. La presenza di strumenti di salvataggio, esportazione e generazione SQL rende il risultato riutilizzabile anche in relazioni, slide o materiali di studio.

Il flusso principale dell'applicazione parte dall'editor ER. Il canvas SVG permette una rappresentazione vettoriale degli elementi e supporta operazioni come selezione, spostamento, zoom, pan, duplicazione, allineamento, annulla e ripeti. Accanto alla vista grafica è presente anche una rappresentazione testuale ERS, sincronizzata con il diagramma, utile per descrivere lo schema in forma più strutturata.

La trasformazione verso lo schema logico è uno degli aspetti più significativi del progetto. Invece di ridurre il passaggio a una conversione automatica non controllata, \builder{} adotta un percorso più esplicito: alcune strutture, come generalizzazioni e attributi composti, vengono trattate in una fase di traduzione ER--ER; successivamente viene costruito il modello logico, con workflow manuale e incrementale. Questo approccio è coerente con un uso didattico, perché rende visibili le decisioni progettuali.

Il repository conferma inoltre la presenza di generazione SQL, esportazione in PNG, SVG e JPEG, salvataggio e caricamento in formato \texttt{.ersp}, reverse engineering SQL e localizzazione in italiano, inglese e albanese. Il progetto si configura quindi come un ambiente completo per progettare, verificare, conservare e condividere schemi di basi di dati.

\chapter{Analisi dei requisiti}

\section{Requisiti funzionali}

I requisiti funzionali individuati riguardano le operazioni che l'applicazione deve mettere a disposizione dell'utente. Il primo requisito è la creazione degli elementi ER principali: entità, relazioni, attributi e collegamenti. A questo si affiancano le operazioni di modifica, cancellazione, rinomina, selezione, spostamento e gestione della disposizione grafica.

Un secondo gruppo di requisiti riguarda la rappresentazione dei vincoli. L'applicazione deve permettere di definire cardinalità, identificatori interni, identificatori esterni, identificatori misti e generalizzazioni ISA. Questi elementi sono fondamentali per costruire schemi concettuali corretti e per preparare la successiva trasformazione logica.

La traduzione verso lo schema logico costituisce un requisito centrale. Il sistema deve produrre tabelle, colonne, chiavi primarie, chiavi esterne e vincoli, gestendo anche casi in cui sono presenti più chiavi candidate o decisioni progettuali alternative. Il repository conferma inoltre la presenza della generazione SQL a partire dal modello logico.

Rientrano tra i requisiti funzionali anche salvataggio e caricamento dei progetti, esportazione in PNG, SVG e JPEG, localizzazione dell'interfaccia e reverse engineering SQL. Quest'ultima funzionalità consente di importare schemi SQL testuali e ottenere una rappresentazione logica o ER, pur con le limitazioni dichiarate nella documentazione tecnica.

\section{Requisiti non funzionali}

I requisiti non funzionali riguardano la qualità complessiva dell'applicazione. Il primo aspetto è l'usabilità: l'utente deve poter comprendere gli strumenti disponibili, interagire con il canvas e ricevere feedback chiari. Un'interfaccia poco leggibile rischierebbe di spostare l'attenzione dai concetti di progettazione ai problemi di utilizzo dello strumento.

La leggibilità grafica è un requisito altrettanto importante. Diagrammi con label sovrapposte, collegamenti poco chiari o spazi vuoti eccessivi negli export risultano meno utili in un contesto didattico. Per questo motivo il progetto presta attenzione a layout, routing, bounding box, gestione degli sfondi e resa degli output.

Altri requisiti riguardano stabilità, compatibilità, responsive, manutenibilità e testabilità. Il formato \texttt{.ersp} deve preservare i progetti nel tempo, anche rispetto a versioni precedenti. La struttura del codice deve rimanere modulare, mentre le funzionalità più delicate devono essere coperte da test. Infine, la chiarezza didattica richiede che le scelte tecniche siano orientate alla comprensione del modello, non soltanto alla presenza di funzioni.

\chapter{Tecnologie utilizzate}

Il progetto è sviluppato con TypeScript, React e Vite. TypeScript consente di definire in modo esplicito le strutture dati dell'applicazione, come nodi, collegamenti, diagrammi, modelli logici, tabelle e file progetto. Questo è particolarmente utile in un'applicazione che manipola dati complessi e deve mantenere coerenza tra rappresentazione grafica, serializzazione e traduzione logica.

React è utilizzato per costruire l'interfaccia utente attraverso componenti. Header, toolbar, canvas, pannelli, modali, loading screen, workspace di traduzione e vista logica sono organizzati in moduli distinti. Vite fornisce l'ambiente di sviluppo, il dev server e la build della single-page application.

Il canvas usa SVG, tecnologia adatta alla rappresentazione di diagrammi vettoriali. Gli stili sono gestiti tramite CSS, mentre la libreria \texttt{lucide-react} supporta il sistema di icone dell'interfaccia. Per la verifica del progetto sono usati \texttt{tsx}, per i test TypeScript, e Playwright, per i test end-to-end nel browser.

\noindent\begin{tabularx}{\textwidth}{>{\bfseries}p{3.5cm}X}
Tecnologia & Ruolo nel progetto \\
\hline
TypeScript & Tipizzazione del dominio applicativo, dei modelli ER/logici e delle funzioni di trasformazione. \\
React & Realizzazione dell'interfaccia utente e dei workspace interattivi. \\
Vite & Dev server, build e preview dell'applicazione web. \\
CSS & Layout, responsive, pannelli, modali, toolbar e stile del canvas. \\
SVG & Rendering vettoriale di diagrammi, nodi, collegamenti e label. \\
npm & Gestione delle dipendenze e degli script di sviluppo, test e build. \\
\texttt{tsx} & Esecuzione dei test TypeScript. \\
Playwright & Test end-to-end dei flussi principali nel browser. \\
Git/GitHub & Versionamento del progetto e organizzazione del repository. \\
GitHub Actions & Automazione della build e del deploy. \\
GitHub Pages & Pubblicazione della build statica. \\
\texttt{lucide-react} & Icone dell'interfaccia. \\
\end{tabularx}

L'insieme di queste tecnologie ha permesso di realizzare un'applicazione interattiva, modulare e verificabile. L'uso di strumenti moderni ha inoltre favorito un'organizzazione più chiara del codice e una maggiore attenzione alla qualità del processo di sviluppo.

\chapter{Architettura e progettazione}

L'architettura di \builder{} è organizzata come applicazione web a componenti. Il file \texttt{src/App.tsx} svolge il ruolo di orchestratore principale: coordina stato del documento, viste, modali, comandi utente, sincronizzazione ERS, undo/redo e passaggio delle proprietà ai diversi workspace. Il repository mostra anche una progressiva separazione di responsabilità verso hook e moduli dedicati.

La cartella \texttt{src/canvas} contiene la logica di rendering e interazione del diagramma ER su SVG. La cartella \texttt{src/components} raccoglie componenti applicativi e condivisi, come header, modali, pannelli, command menu, loading screen e preview SQL reverse. Le cartelle \texttt{src/inspector} e \texttt{src/toolbar} gestiscono rispettivamente la modifica delle proprietà degli elementi e gli strumenti disponibili nel canvas.

Le cartelle \texttt{src/translation} e \texttt{src/logical} sono dedicate al percorso di trasformazione e alla gestione dello schema logico. La localizzazione è concentrata in \texttt{src/i18n}, mentre \texttt{src/types} definisce i contratti TypeScript condivisi. La cartella \texttt{src/utils} contiene la logica pura e testabile: parser, serializzazione, validazione, layout, export, traduzione logica, SQL reverse e gestione dei file progetto.

\begin{verbatim}
Interfaccia utente
        v
Gestione dello stato
        v
Modello ER
        v
Traduzione logica
        v
Esportazione / salvataggio
\end{verbatim}

Questa organizzazione rende più chiara la separazione tra interfaccia, stato, modello dati, traduzione ed export. La centralità dei tipi TypeScript permette di descrivere il dominio applicativo in modo coerente, mentre la presenza di utility pure facilita la scrittura dei test. La compatibilità dei file salvati e il ripristino della sessione sono aspetti progettuali importanti, perché consentono all'utente di proseguire il lavoro nel tempo.

\chapter{Funzionalità sviluppate e consolidate}

\section{Editor ER}

L'editor ER rappresenta il nucleo dell'applicazione. L'utente può creare entità, relazioni e attributi, collegarli tra loro, indicare cardinalità e modellare vincoli. Il canvas supporta selezione, spostamento, zoom, pan, duplicazione, allineamento e undo/redo. Queste funzionalità rendono l'ambiente adatto alla costruzione progressiva di schemi anche complessi.

Particolare attenzione è stata posta alla gestione grafica degli elementi. La disposizione degli attributi, il posizionamento delle label e il routing dei collegamenti incidono direttamente sulla leggibilità dello schema. Il progetto include quindi logiche dedicate a layout, collision avoidance e stabilità della geometria.

Gli identificatori interni, esterni e misti costituiscono una parte rilevante del lavoro. La loro corretta gestione è fondamentale perché incide sia sulla validità dello schema ER sia sulla successiva traduzione verso il modello logico. Anche le generalizzazioni ISA sono supportate con vincoli specifici, come total/partial e disjoint/overlap.

\section{Traduzione verso lo schema logico}

La traduzione verso lo schema logico è organizzata come processo guidato. Il sistema non si limita a convertire immediatamente il diagramma, ma prevede una pipeline ER--ER per risolvere strutture come generalizzazioni e attributi composti. Successivamente, il workspace logico permette di costruire il modello relazionale attraverso decisioni più esplicite.

Il modello logico comprende tabelle, colonne, chiavi primarie, chiavi esterne e vincoli. In presenza di più identificatori candidati, l'applicazione prevede una scelta esplicita della chiave primaria, rendendo visibile una decisione progettuale importante. Questo comportamento è particolarmente utile in un contesto didattico, perché aiuta a comprendere che la traduzione non è sempre un processo meccanico.

\section{Esportazione, salvataggio e caricamento}

\builder{} supporta l'esportazione in PNG, SVG e JPEG. Il changelog documenta miglioramenti specifici relativi a bounding box reale dello schema, riduzione dello spazio vuoto, sfondo trasparente per PNG e SVG, sfondo bianco per JPEG e conservazione degli stili nello SVG standalone. Queste caratteristiche rendono gli output più adatti all'uso in documenti e presentazioni.

Il salvataggio avviene tramite formato \texttt{.ersp}, che conserva diagramma ER, workspace di traduzione, workspace logico, vista corrente, viewport e metadati. La compatibilità con file legacy e il ripristino della sessione mostrano attenzione alla continuità del lavoro dell'utente.

\section{Interfaccia e localizzazione}

L'interfaccia include toolbar, pannelli, modali, command menu, loading screen e scorciatoie da tastiera. Il lavoro sul responsive ha riguardato desktop, tablet e telefono, con l'obiettivo di ridurre sovrapposizioni e rendere i controlli più accessibili.

La localizzazione in italiano, inglese e albanese è integrata nell'applicazione tramite dizionari e selettore lingua. Questa funzionalità amplia l'accessibilità dello strumento e rafforza il suo possibile uso didattico in contesti diversi.

\chapter{Test, verifica e miglioramenti}

La qualità di \builder{} è stata verificata tramite test unitari, test di integrazione e test end-to-end. I test nella cartella \texttt{test/} coprono parser ERS, serializzazione, traduzione ER, traduzione logica, identificatori, layout, export, file progetto, SQL reverse, localizzazione e sessione. I test Playwright nella cartella \texttt{tests/e2e/} verificano flussi reali nel browser, come loading screen e cambio lingua.

\noindent\begin{tabularx}{\textwidth}{>{\bfseries}p{4cm}p{4cm}X}
Area testata & Tipo di verifica & Scopo \\
\hline
Parsing e serializzazione ERS & Test unitari e di integrazione & Verificare la coerenza tra rappresentazione testuale e diagramma. \\
Formato \texttt{.ersp} & Test di integrazione & Controllare salvataggio, caricamento, compatibilità e migrazione legacy. \\
Traduzione ER--ER & Test unitari & Verificare generalizzazioni, attributi composti e decisioni di traduzione. \\
Vista logica & Test unitari e di integrazione & Controllare tabelle, chiavi, foreign key, vincoli e workflow manuale. \\
Export & Test unitari & Verificare PNG, SVG, JPEG, sfondi, bounds e stili. \\
Layout & Test unitari & Evitare sovrapposizioni di attributi, cardinalità, edge label e label FK. \\
SQL reverse & Test unitari e di integrazione & Controllare parser SQL, modello logico, diagramma ER e layout generato. \\
Localizzazione & Test unitari ed E2E & Verificare dizionari, selettore lingua e testi dell'interfaccia. \\
Sessione workspace & Test unitari & Controllare ripristino, sanitizzazione e compatibilità dei dati salvati. \\
\end{tabularx}

I miglioramenti introdotti hanno riguardato in particolare layout, label, export, sfondi, compatibilità dei salvataggi, responsive, localizzazione e stabilità. In diversi casi, la correzione di un problema è stata accompagnata dall'aggiunta di test, in modo da ridurre il rischio di regressioni.

Dal punto di vista del tirocinio, questa attività ha permesso di comprendere meglio l'importanza della verifica. In un progetto grafico e interattivo, testare non significa soltanto controllare il risultato visibile, ma anche verificare parser, trasformazioni, serializzazione e regole di dominio.

\chapter{Supervisione della tutor e competenze acquisite}

La Prof.ssa Alessandra Lumini ha svolto il ruolo di tutor del tirocinio, supervisionando il percorso e fornendo indicazioni utili per orientare il lavoro. Il suo contributo ha aiutato a mantenere l'attenzione non solo sulla realizzazione tecnica, ma anche sulla chiarezza concettuale, sulla correttezza degli schemi, sulla leggibilità grafica e sull'utilità didattica dello strumento.

Le indicazioni della tutor hanno contribuito a definire le priorità del progetto. In particolare, hanno favorito un'attenzione costante alla coerenza tra schema ER e schema logico, alla semplicità dell'interfaccia e alla necessità di mantenere il codice organizzato e verificabile. Non si tratta di citazioni dirette, ma di una sintesi prudente del ruolo di supervisione svolto durante il tirocinio.

\noindent\begin{tabularx}{\textwidth}{>{\bfseries}p{4cm}X}
Aspetto supervisionato & Effetto sul lavoro \\
\hline
Chiarezza concettuale & Maggiore attenzione al significato degli elementi ER e alla loro trasformazione logica. \\
Correttezza degli schemi & Controllo di entità, relazioni, cardinalità, identificatori e generalizzazioni. \\
Leggibilità grafica & Cura di layout, label, collegamenti e output esportati. \\
Usabilità & Miglioramento di toolbar, pannelli, modali, flussi utente e responsive. \\
Valore didattico & Collegamento tra funzionalità tecniche e comprensione della progettazione di basi di dati. \\
Manutenibilità & Separazione tra interfaccia e logica di dominio, uso di tipi condivisi e test. \\
Revisione progressiva & Consolidamento incrementale del progetto e maggiore attenzione alla qualità complessiva. \\
\end{tabularx}

\section{Competenze tecniche}

Il tirocinio ha permesso di consolidare competenze nello sviluppo web con React, TypeScript e Vite, nella gestione di canvas SVG e nella modellazione di strutture dati per schemi ER e logici. Sono state inoltre rafforzate competenze su esportazione immagini, salvataggio, serializzazione, generazione SQL, test software, Git e debugging.

\section{Competenze metodologiche e trasversali}

Sul piano metodologico, il lavoro ha richiesto autonomia, problem solving, pianificazione e revisione progressiva. Il confronto con la tutor ha favorito la capacità di ricevere feedback, trasformarlo in priorità operative e comunicare in modo più chiaro le scelte progettuali. Sono state inoltre sviluppate competenze nella documentazione del lavoro e nell'attenzione all'utente finale.

\chapter{Conclusioni e sviluppi futuri}

Il tirocinio ha portato alla preparazione e al consolidamento di un progetto web articolato, capace di supportare diverse fasi della progettazione di basi di dati. \builder{} integra editor ER, vista logica, esportazione, salvataggio, localizzazione, test e funzionalità di reverse engineering SQL. Questi risultati mostrano un'evoluzione significativa dello strumento e ne confermano il possibile valore didattico.

L'esperienza ha avuto anche un valore formativo rilevante. Il progetto ha richiesto di lavorare su un'applicazione reale, con molte parti interdipendenti, e di collegare le scelte tecniche agli obiettivi di chiarezza, correttezza e usabilità. La supervisione della Prof.ssa Alessandra Lumini ha contribuito a mantenere questo orientamento, aiutando a considerare il software non solo come prodotto tecnico, ma come strumento di apprendimento.

Restano alcuni limiti realistici. Il layout automatico può essere ulteriormente migliorato per schemi molto grandi o densi. Il parser SQL, pur utile, non mira a coprire l'intera grammatica dei diversi dialetti. L'accessibilità e la documentazione utente possono essere ampliate, così come i test end-to-end. Inoltre, funzionalità collaborative online non risultano presenti e potrebbero rappresentare un'evoluzione futura.

Tra gli sviluppi futuri si possono indicare l'introduzione di esempi didattici integrati, una guida utente più completa, il miglioramento dell'accessibilità, l'ampliamento dei costrutti EER, il rafforzamento del parser SQL, ulteriori test e un layout automatico ancora più robusto. Un possibile sviluppo più avanzato potrebbe riguardare la collaborazione online, utile in scenari di revisione o esercitazione condivisa.

In conclusione, il tirocinio ha permesso di consolidare competenze tecniche e metodologiche, affrontando un progetto significativo per la progettazione di basi di dati. Il lavoro svolto ha richiesto attenzione alla qualità del codice, alla correttezza dei modelli, alla chiarezza dell'interfaccia e alla verifica delle funzionalità. Per questo motivo, l'esperienza rappresenta un passaggio importante nel percorso formativo universitario e nella maturazione di un approccio più consapevole allo sviluppo software.

\appendix
\chapter{Dati da completare prima della consegna}

\begin{itemize}
    \item \placeholder{Nome e cognome studente}.
    \item \placeholder{Matricola da completare}.
    \item \placeholder{Corso di laurea da completare}.
    \item \placeholder{Anno accademico da completare}.
    \item \placeholder{Periodo del tirocinio da completare}.
    \item \placeholder{Ore/CFU da completare}.
    \item \placeholder{Struttura ospitante da completare}.
    \item \placeholder{Link repository o sito pubblicato da completare}.
    \item Eventuali screenshot dell'interfaccia, dello schema ER, della vista logica e dell'export.
\end{itemize}

\cleardoublepage

\chapter*{Bibliografia e sitografia}
\addcontentsline{toc}{chapter}{Bibliografia e sitografia}

\begin{itemize}
    \item Peter P. Chen, \textit{The Entity-Relationship Model: Toward a Unified View of Data}, ACM Transactions on Database Systems, 1976.
    \item Ramez Elmasri, Shamkant B. Navathe, \textit{Fundamentals of Database Systems}, Pearson.
    \item Documentazione ufficiale React: \url{https://react.dev/}.
    \item Documentazione ufficiale TypeScript: \url{https://www.typescriptlang.org/docs/}.
    \item Documentazione ufficiale Vite: \url{https://vite.dev/guide/}.
    \item Documentazione ufficiale Playwright: \url{https://playwright.dev/}.
    \item Repository o sito del progetto \builder: \placeholder{Link da completare}.
\end{itemize}

\end{document}
